% =====================================================================
%  ΚΕΦΑΛΑΙΟ 1 — ΕΙΣΑΓΩΓΗ
%  Κάθε αριθμητικό στοιχείο του κεφαλαίου συνοδεύεται από πηγή και από
%  το πειραματικό περιβάλλον στο οποίο παρήχθη. Καμία μέτρηση της
%  παρούσας εργασίας δεν εμφανίζεται σε αυτό το κεφάλαιο.
%
%  Η ΑΓΓΛΙΚΗ ΕΙΝΑΙ ΤΟ ΠΡΩΤΟΤΥΠΟ. Αυτό το αρχείο ακολουθεί το
%  content/en/ch1.tex, μπλοκ προς μπλοκ. Δες docs/SYNC.md.
% =====================================================================

%%blk intro.head from=91d794864574
\chapter{Εισαγωγή: το πρόβλημα της συναίνεσης και το ερώτημα της εργασίας}
\label{ch:intro}

%%blk intro.head.roadmap from=b89a4362725b
Το κεφάλαιο κάνει τέσσερα πράγματα, με αυτή τη σειρά. Διατυπώνει το
πρόβλημα που λύνει η συναίνεση και τα δύο εμπόδια που το κάνουν δύσκολο,
τα σφάλματα και τις καθυστερήσεις. Εξηγεί γιατί δεν το έλυσε οριστικά
ένα πρωτόκολλο, και γιατί οι δημοσιευμένες συγκρίσεις μεταξύ πρωτοκόλλων
δεν μεταφέρονται από τη μία εργασία στην άλλη. Ύστερα δηλώνει τον σκοπό
της εργασίας, τα ερευνητικά ερωτήματα και τα όριά της. Κλείνει με τον
κανόνα που ακολουθεί η εργασία για την ποιότητα των πηγών της, γιατί μια
σύγκριση αξίζει ακριβώς όσο και τα τεκμήρια από κάτω της.

%%blk intro.problem.head from=e3c836bd5a2a
\section{Συμφωνία υπό σφάλματα και καθυστερήσεις}
\label{sec:problem}

%%blk intro.problem.01 from=e1271f592a84
Ένα κατανεμημένο σύστημα αποτελείται από διεργασίες χωρικά
διαχωρισμένες, που επικοινωνούν αποκλειστικά με ανταλλαγή μηνυμάτων. Ένα
σύστημα είναι κατανεμημένο όταν η καθυστέρηση μετάδοσης δεν είναι
αμελητέα σε σχέση με τον χρόνο που μεσολαβεί ανάμεσα σε δύο γεγονότα της
ίδιας διεργασίας \cite{lamport1978time}. Ο ορισμός αυτός έχει μία
συνέπεια: δεν υπάρχει καθολικό ρολόι. Για δύο γεγονότα σε διαφορετικές
διεργασίες είναι ενίοτε αδύνατο να ειπωθεί ποιο συνέβη πρώτο. Η μόνη
εγγενής σχέση διάταξης είναι η μερική σχέση \en{happened-before}
\cite{lamport1978time}. Κάθε υπηρεσία που οφείλει να συμπεριφέρεται σαν
ένα ενιαίο, αξιόπιστο μηχάνημα πρέπει επομένως να κατασκευάσει η ίδια τη
σειρά αυτή, επειδή το δίκτυο δεν της τη δίνει.

%%blk intro.problem.02 from=63ea30011c9d
Η κανονική αφαίρεση με την οποία κατασκευάζεται αυτή η σειρά είναι η
\emph{αναπαραγωγή μηχανής καταστάσεων} (\en{State Machine Replication},
SMR). Η SMR μοντελοποιεί την υπηρεσία ως ντετερμινιστική μηχανή
καταστάσεων και την αναπαράγει σε $n$ κόμβους. Αν όλοι οι σωστοί κόμβοι
εκτελέσουν την ίδια ακολουθία εντολών με την ίδια σειρά, θα καταλήξουν
στην ίδια κατάσταση και θα παραγάγουν τις ίδιες εξόδους
\cite{schneider1990implementing}. Το πρόβλημα της αξιοπιστίας ανάγεται
έτσι σε ένα και μόνο πρόβλημα συντονισμού, τη \emph{συναίνεση}
(\en{consensus}): τη συμφωνία των σωστών κόμβων σε μία κοινή τιμή για
κάθε θέση της ακολουθίας. Το κριτήριο ορθότητας που περιμένει ο πελάτης
από την αναπαραγόμενη υπηρεσία είναι η \emph{γραμμικοποιησιμότητα}
(\en{linearizability}). Κάθε πράξη φαίνεται να συμβαίνει ακαριαία σε
κάποιο σημείο ανάμεσα στην κλήση της (\en{invocation}) και στην απόκρισή
της, και οι πράξεις που δεν επικαλύπτονται χρονικά διατηρούν τη σειρά
τους \cite{herlihy1990linearizability}.

%%blk intro.problem.03 from=3d8d7e1d6b7d
Ένα πρωτόκολλο συναίνεσης κρίνεται από δύο ιδιότητες. Η \emph{ασφάλεια}
(\en{safety}) απαιτεί να μην αποφασίσουν ποτέ δύο σωστοί κόμβοι
διαφορετικά για την ίδια θέση. Η \emph{ζωτικότητα} (\en{liveness})
απαιτεί κάθε αίτημα σωστού πελάτη να αποφασίζεται τελικά. Το θεμελιώδες
εμπόδιο είναι ότι κανένα ντετερμινιστικό πρωτόκολλο δεν εξασφαλίζει τις
δύο ιδιότητες ταυτόχρονα σε πλήρως ασύγχρονο σύστημα. Οι Fischer, Lynch
και Paterson απέδειξαν ότι \emph{κάθε ντετερμινιστικό} πρωτόκολλο για το
πρόβλημα αυτό έχει τη δυνατότητα μη τερματισμού, ακόμη και με μία μόνο
ελαττωματική διεργασία \cite{fischer1985impossibility}. Το αποτέλεσμα είναι δομικό και
όχι τεχνικός περιορισμός: σε ασύγχρονο σύστημα καμία διεργασία δεν μπορεί
να διακρίνει μια νεκρή διεργασία από μια αργή.

%%blk intro.problem.04 from=5f7e2e3ff158
Όλα τα πρωτόκολλα που εξετάζει η εργασία χρησιμοποιούν την ίδια διέξοδο,
το μοντέλο της \emph{μερικής συγχρονίας} (\en{partial synchrony})
\cite{dwork1988consensus}. Στη μορφή που υιοθετούν τα σύγχρονα
πρωτόκολλα, υπάρχει άγνωστη χρονική στιγμή. Είναι η GST, το \en{Global
Stabilization Time}. Μετά την GST, το δίκτυο παραδίδει τα μηνύματα μεταξύ
σωστών κόμβων εντός γνωστού ορίου $\Delta$. Έτσι ακριβώς διατυπώνουν το
μοντέλο τα βυζαντινά πρωτόκολλα της εργασίας \cite{yin2019hotstuff} και
\cite{malkhi2023hotstuff2}\prep. Η ασφάλεια διατηρείται πάντοτε, ακόμη
και πριν την GST, και η ζωτικότητα εξασφαλίζεται μόνο μετά από αυτήν.

%%blk intro.problem.05 from=8a0fcb4b9434
Η συνέπεια για κάθε πειραματική αποτίμηση είναι σημαντική. Ο χρόνος
ανάκαμψης μετά την απώλεια του ηγέτη εξαρτάται σε μεγάλο βαθμό από τη
ρύθμιση των χρονομέτρων, όχι μόνο από τον αλγόριθμο. Δύο εργασίες το
διατυπώνουν ρητά. Η μελέτη αναπαραγωγιμότητας του Raft το σημειώνει
πρώτη \cite{howard2015raftrefloated}. Οι Berger κ.ά.\ το επαναλαμβάνουν
για βυζαντινά συστήματα. Γράφουν ότι ο χρόνος ανάκαμψης εξαρτάται από τις
παραμέτρους των χρονομέτρων και ότι \emph{δεν αποκλείουν} πως στενότερες
ρυθμίσεις θα ήταν εφικτές \cite{berger2023simulation}.

%%blk intro.problem.06 from=1424ac3bc7e6
Η δεύτερη διάσταση που οργανώνει ολόκληρη τη βιβλιογραφία είναι το
\emph{μοντέλο σφαλμάτων}. Στο μοντέλο κατάρρευσης (\en{crash-fault}) ένας
κόμβος μπορεί μόνο να σταματήσει. Στο βυζαντινό μοντέλο ένας κόμβος
μπορεί να συμπεριφερθεί αυθαίρετα, να στείλει αντιφατικά μηνύματα σε
διαφορετικούς παραλήπτες και να συντονιστεί με άλλους ελαττωματικούς
κόμβους \cite{lamport1982byzantine}. Η διάκριση δεν είναι φιλοσοφική.
Καθορίζει πόσοι κόμβοι χρειάζονται, όπως δείχνει η
Ενότητα~\ref{sec:notonce}.

%%blk intro.notonce.head from=cd630a475d44
\section{Γιατί το πρόβλημα δεν λύθηκε μία φορά}
\label{sec:notonce}

%%blk intro.notonce.01 from=e1dc502bef0c
Το Viewstamped Replication δημοσιεύτηκε το 1988. Είναι το πρώτο
πρωτόκολλο της γραμμής που εξετάζει η εργασία. Από τότε έχουν περάσει
σχεδόν τέσσερις δεκαετίες, και η παραγωγή νέων πρωτοκόλλων δεν έχει
σταματήσει. Το Σχήμα~\ref{fig:timeline} δίνει τη χρονολογική σειρά με
έναν ρητό κανόνα: κάθε πρωτόκολλο τοποθετείται στην \emph{πρώτη
δημοσίευσή του με κρίση}, σε πρακτικά συνεδρίου ή σε περιοδικό. Τρία από
τα δέκα πρωτόκολλα δεν έχουν τέτοια δημοσίευση. Το σχήμα τα σχεδιάζει με
κενό δείκτη, τα τοποθετεί στην πρώτη δημόσια διάθεση του σχεδιασμού τους
και σημαίνει την πηγή τους. Η εργασία δηλώνει τον κανόνα εδώ, επειδή η
λέξη «χρονολογικά» χωρίς κανόνα γίνεται αυθαίρετη, και η αυθαιρεσία
φαίνεται.

%%blk intro.notonce.fig.timeline from=055513a6b49d
\dgfigure{fig:timeline}%
{Χρονολογική σειρά των πρωτοκόλλων που εξετάζει η εργασία, με κανόνα
τοποθέτησης την πρώτη δημοσίευση με κρίση.}%
{Χρονολογική σειρά των πρωτοκόλλων που εξετάζει η εργασία, με κανόνα
τοποθέτησης την πρώτη δημοσίευση με κρίση. \textbf{Γεμάτος δείκτης} =
πρώτη δημοσίευση με κρίση, η οποία και ορίζει τη θέση: VR
\cite{oki1988viewstamped}, Paxos \cite{lamport1998parttime}, PBFT
\cite{castro1999pbft}, ZAB \cite{junqueira2011zab}, EPaxos
\cite{moraru2013epaxos}, Raft \cite{ongaro2014raft}, HotStuff
\cite{yin2019hotstuff}. \textbf{Κενός δείκτης} = δεν υπάρχει δημοσίευση
με κρίση. Η θέση δίνεται τότε από την πρώτη δημόσια διάθεση του
σχεδιασμού, και το σχήμα τη σημαίνει. KRaft: KIP-500, τεχνική
τεκμηρίωση\grey\ του 2019. Η γενική διάθεση ήρθε με το Kafka~3.3 το 2022
\cite{kafka2019kip500}. Tendermint: αλγοριθμικό \en{preprint}\prep\ του
2018 \cite{buchman2018gossip}, \en{whitepaper}\grey\ του 2014
\cite{kwon2014tendermint}, και πρώτη κριθείσα εργασία με μετρήσεις του
συστήματος το 2021 \cite{cason2021tendermint}\unver. HotStuff-2: εκτεταμένη
περίληψη σε ePrint\prep\ \cite{malkhi2023hotstuff2}. \emph{Δύο
επιφυλάξεις προέλευσης.} Το \cite{lamport1998parttime} φέρει ημερομηνία
παραλαβής τον Ιανουάριο του 1990 και αποδοχής τον Μάρτιο του 1998, αλλά
τα δύο αντίτυπα του \en{corpus} τη \emph{διατυπώνουν} διαφορετικά (βλ.\
\S\ref{sec:sources}). Η τεχνική αναφορά SRC~RR-49 χρονολογείται ευρέως
στο 1989. Δεν την εντοπίσαμε ως τεκμήριο και γι' αυτό δεν ορίζει τη θέση.
Το PBFT έχει και εκτενή εκδοχή σε περιοδικό το 2002
\cite{castro2002pbftrecovery}. Θεωρητικά ορόσημα:
\cite{lamport1978time,lamport1982byzantine,fischer1985impossibility,%
dwork1988consensus,schneider1990implementing,herlihy1990linearizability}.}%
{% fig-timeline — first-refereed-publication rule, CFT above the axis, BFT below.
% Structure figure: dgflow. The only geometry it sets is the year-to-length
% mapping, and that is tied to \dgwidth so the timeline is exactly as wide as
% every other figure in the thesis: 51 years (1977-2028) across \dgwidth.
% Every human-readable string is a label key; protocol names, venues and years
% are identical in both editions and stay literal.
\begin{dgflow}[x={\dgwidth/51}]
  \def\yzero{1977}

  % Theoretical milestones. A third category — neither CFT nor BFT protocols —
  % so dgAux rather than either series colour.
  \node[dgnote, text=dgAux, anchor=west] at (0,4.45)
    {\textbf{\fl{timeline.milestones.title}}\\
     1978 \fl{timeline.milestone.clocks}\, ·\, 1982 \fl{timeline.milestone.byzmodel}\, ·\, 1985 FLP\\
     1988 \fl{timeline.milestone.partialsync}\, ·\, 1990 SMR\, ·\, 1990 \fl{timeline.milestone.linearizability}};

  % The time axis and its five-year ticks. The year LABELS are drawn last, so
  % that their dgtag background sits over the leader lines rather than under.
  \draw[dgaxis] (0,0) -- (51,0);
  \foreach \yr in {1980,1985,1990,1995,2000,2005,2010,2015,2020,2025}{
    \pgfmathsetmacro{\xx}{\yr-\yzero}
    \draw[dgaxis,-] (\xx,-0.08) -- (\xx,0.08);}

  % Which half is which. At the left end, where the axis is still empty.
  \node[dgtag, fill=dgPaper, text=dgCFT, anchor=west] at (0.4, 1.00) {\textbf{CFT}};
  \node[dgtag, fill=dgPaper, text=dgBFT, anchor=west] at (0.4,-1.00) {\textbf{BFT}};

  % Crash-fault protocols, above the axis. Three stagger levels: 2011-2019 is
  % 1.7 cm of axis carrying 5.6 cm of label, and the anchors are chosen so that
  % no leader line ever runs through another protocol's label.
  % year / level / anchor / filled marker? / label
  \foreach \yr/\lvl/\anc/\solid/\lbl in {%
      1988/1.00/{south}/{1}/{VR\\PODC '88},%
      1998/1.00/{south}/{1}/{Paxos\\TOCS '98},%
      2011/1.00/{south east}/{1}/{ZAB\\DSN '11},%
      2013/1.95/{south east}/{1}/{EPaxos\\SOSP '13},%
      2014/2.90/{south}/{1}/{Raft\\ATC '14},%
      2019/1.00/{south west}/{0}/{KRaft\\KIP-500 '19\grey}}{
    \pgfmathsetmacro{\xx}{\yr-\yzero}
    \draw[dgarrow,-,draw=dgCFT] (\xx,0.07) -- (\xx,\lvl);
    \ifnum\solid=1 \fill[dgCFT] (\xx,0.07) circle (1.7pt);
    \else \draw[dgarrow,-,draw=dgCFT,fill=dgPaper] (\xx,0.07) circle (1.7pt);\fi
    \node[dgtag, fill=dgPaper, text=dgCFT, anchor=\anc] at (\xx,\lvl) {\lbl};}

  % Byzantine-fault protocols, below the axis. Same three levels, same reason:
  % 2018, 2019 and 2023 are 1.1 cm apart and their labels are 1.5 cm wide.
  \foreach \yr/\lvl/\anc/\solid/\lbl in {%
      1999/-1.00/{north}/{1}/{PBFT\\OSDI '99},%
      2018/-1.00/{north east}/{0}/{Tendermint\\arXiv '18\prep},%
      2019/-1.95/{north east}/{1}/{HotStuff\\PODC '19},%
      2023/-2.90/{north}/{0}/{HotStuff-2\\ePrint '23\prep}}{
    \pgfmathsetmacro{\xx}{\yr-\yzero}
    \draw[dgarrow,-,draw=dgBFT] (\xx,-0.07) -- (\xx,\lvl);
    \ifnum\solid=1 \fill[dgBFT] (\xx,-0.07) circle (1.7pt);
    \else \draw[dgarrow,-,draw=dgBFT,fill=dgPaper] (\xx,-0.07) circle (1.7pt);\fi
    \node[dgtag, fill=dgPaper, text=dgBFT, anchor=\anc] at (\xx,\lvl) {\lbl};}

  % Year labels last: 1988, 1999, 2011, 2014 and 2019 all sit within half a
  % label of a five-year tick, so something has to be on top and it should be
  % the year.
  \foreach \yr in {1980,1985,1990,1995,2000,2005,2010,2015,2020,2025}{
    \pgfmathsetmacro{\xx}{\yr-\yzero}
    \node[dgtag, fill=dgPaper, text=dgGrey, anchor=north] at (\xx,-0.16) {\yr};}
\end{dgflow}
}

%%blk intro.notonce.02 from=440b63533c93
Η επανάληψη δεν οφείλεται σε αδυναμία των προηγούμενων λύσεων να είναι
ορθές. Οφείλεται στο ότι η ορθότητα είναι μόνο η μία διάσταση. Κάθε
πρωτόκολλο είναι μια \emph{θέση} σε έναν χώρο σχεδίασης, και κάθε
βελτίωση εκεί αγοράζεται με συγκεκριμένο τίμημα. Η πρόσφατη ενοποιητική
εργασία των Amiri κ.ά.\ το διατυπώνει ρητά. Περιγράφει τα βυζαντινά
πρωτόκολλα ως σημεία ενός χώρου με περισσότερες από δέκα διαστάσεις
σχεδίασης. Στον χώρο αυτόν κάθε επιλογή είναι μια αμφιμονοσήμαντη
συνάρτηση από πρωτόκολλο σε πρωτόκολλο \cite{amiri2024bedrock}.

%%blk intro.notonce.03 from=11de5b63b6e0
Η σαφέστερη ένδειξη αυτού είναι η αριθμητική των \en{quorums}
(Σχήμα~\ref{fig:quorum}). Το πλήθος των κόμβων δεν αυξάνεται αυθαίρετα.
Μετακινείται μόνο όταν αγοράζει κάτι συγκεκριμένο και κατονομάσιμο. Με
σφάλματα κατάρρευσης αρκεί ένα πλειοψηφικό \en{quorum}. Μια συστάδα
$n = 2f+1$ κόμβων παραμένει διαθέσιμη όσο λειτουργεί μια πλειοψηφία, άρα
ανέχεται $f$ σφάλματα \cite{ongaro2014raft}. Η γενική διατύπωση της
ασφάλειας ως τομή πλειοψηφικών συνόλων ανάγεται στο
\cite{lamport1998parttime}.

%%blk intro.notonce.04 from=cb4943a00b1a
Με βυζαντινά σφάλματα και μη αυθεντικοποιημένα μηνύματα απαιτούνται
τουλάχιστον $n = 3f+1$ κόμβοι. Το φράγμα διατυπώνεται πρώτη φορά το 1980,
ως πρόβλημα διαδραστικής συνέπειας \cite{pease1980reaching}. Η αλληγορία
των βυζαντινών στρατηγών είναι μεταγενέστερη \cite{lamport1982byzantine}.
Το ίδιο φράγμα ισχύει και για τα πρωτόκολλα μερικής συγχρονίας που
εξετάζονται εδώ \cite{dwork1988consensus,amiri2024bedrock}. Εκεί, όμως,
δεν προκύπτει από την απουσία υπογραφών, και η
Ενότητα~\ref{sec:bounds} δείχνει γιατί.

%%blk intro.notonce.05 from=c7c6ac12e911
Το έμπιστο υλικό (\en{trusted execution environment}) εμποδίζει έναν
ελαττωματικό κόμβο να στείλει αντιφατικά μηνύματα χωρίς να γίνει
αντιληπτός. Αν προστεθεί τέτοιο υλικό, το φράγμα επιστρέφει στο $2f+1$
\cite{amiri2024bedrock}. Η εξοικονόμηση μιας φάσης επικοινωνίας κόστιζε
αρχικά $5f+1$ κόμβους \cite{martin2006fab}\unver. Με αυθεντικοποίηση αρκούν
$5f-1$. Ο αριθμός αυτός είναι και το αποδεδειγμένο κάτω φράγμα για
βυζαντινή συναίνεση δύο βημάτων \cite{kuznetsov2021optimal}\unver, και
επιτυγχάνεται στην πράξη \cite{amiri2024bedrock}. Τέλος, η υποστήριξη
δίκαιης διάταξης (\en{order-fairness}) απαιτεί τουλάχιστον $4f+1$ κόμβους
στη \emph{λιγότερο} απαιτητική παραμετροποίησή της. Καθώς η παράμετρος
$\gamma$ πλησιάζει το $0{,}5$ και η απαίτηση δικαιοσύνης γίνεται
αυστηρότερη, το φράγμα $n > 4f/(2\gamma-1)$ αυξάνεται απεριόριστα
\cite{kelkar2020orderfairness,amiri2024bedrock}\unver.

%%blk intro.notonce.fig.quorum from=e304361d83af
\dgfigure{fig:quorum}%
{Αριθμητική \protect\en{quorum}: το απαιτούμενο πλήθος κόμβων $n$ ως
συνάρτηση των ανεκτών σφαλμάτων $f$.}%
{Αριθμητική \protect\en{quorum}: το απαιτούμενο πλήθος κόμβων $n$ ως
συνάρτηση των ανεκτών σφαλμάτων $f$. Δεν πρόκειται για μετρήσεις. Είναι
καθαρή αριθμητική από φράγματα που δηλώνονται στη βιβλιογραφία. Για
$f=3$: η βυζαντινή ανοχή κοστίζει 10 κόμβους έναντι 7, η δίκαιη διάταξη
13 και η εξοικονόμηση μίας φάσης 16. \emph{Πηγές.} Το
\cite{amiri2024bedrock} συγκεντρώνει και τα πέντε φράγματα: $3f{+}1$,
$2f{+}1$ και $5f{\pm}1$ στο \S3.2 E1, έμπιστο υλικό στο Παράρτ.~A, O7,
και δίκαιη διάταξη στο \S3.3 Q1 και στη σχεδιαστική επιλογή 13. Για τα
φράγματα αυτά είναι δευτερογενής πηγή. Τα αποδίδει στα
\cite{pease1980reaching,lamport1982byzantine} για το $3f{+}1$ (η
πρωτογενής διατύπωση είναι του 1980, βλ.\ Ενότητα~\ref{sec:bounds}), στο
\cite{martin2006fab}\unver\ για το $5f{+}1$, στα \cite{kuznetsov2021optimal}\unver\ και
\cite{abraham2021goodcase}\unver\ για το $5f{-}1$ (το \cite{amiri2024bedrock}
παραθέτει και τα δύο), και στο \cite{kelkar2020orderfairness}\unver\ για το
$4f{+}1$. Από αυτές τις πηγές, μόνο τα \cite{pease1980reaching} και
\cite{lamport1982byzantine} επαληθεύτηκαν ως προς τα βιβλιογραφικά τους
στοιχεία, μέσω ευρετηρίου. Τα υπόλοιπα τέσσερα δεν επαληθεύτηκαν και τα
στοιχεία τους προέρχονται από τον κατάλογο αναφορών του
\cite{amiri2024bedrock}. \emph{Κανένα από τα έξι δεν βρίσκεται στο
\en{corpus}.} \emph{Σημείωση ακρίβειας.} Το \cite{amiri2024bedrock}
διακρίνει το μέγεθος του \en{quorum} από το μέγεθος του δικτύου. Σε
σχεδιάσεις με ενεργούς και παθητικούς κόμβους το \en{quorum} πέφτει στο
$2f{+}1$ ενώ το $n$ παραμένει $3f{+}1$.}%
{% fig-quorum — required n as a function of tolerated faults f. Pure arithmetic
% from bounds stated in the literature; no measurements.
% Data figure: dgdata. The five series are plotted in the order that makes the
% standard cycle list hand them their meanings: dgCFT crash, dgBFT Byzantine,
% then dgAux / dgSand / dgPurple for the three further bounds.
\begin{dgdata}[
  xlabel={\fl{quorum.xlabel} $f$},
  ylabel={\fl{quorum.ylabel} $n$},
  xmin=0.7, xmax=5.3, ymin=0, ymax=28,
  xtick={1,2,3,4,5}, ytick={0,5,10,15,20,25},
  legend style={at={(0.5,-0.30)}, anchor=north}, legend columns=1,
]
\addplot coordinates {(1,3)(2,5)(3,7)(4,9)(5,11)};
  \addlegendentry{$2f{+}1$ — \fl{quorum.crash}}
\addplot coordinates {(1,4)(2,7)(3,10)(4,13)(5,16)};
  \addlegendentry{$3f{+}1$ — \fl{quorum.byzantine}}
\addplot coordinates {(1,5)(2,9)(3,13)(4,17)(5,21)};
  \addlegendentry{$4f{+}1$ — \fl{quorum.fairorder} ($\gamma=1$)}
\addplot coordinates {(1,6)(2,11)(3,16)(4,21)(5,26)};
  \addlegendentry{$5f{+}1$ — \fl{quorum.onephaseless} (FaB)}
\addplot coordinates {(1,4)(2,9)(3,14)(4,19)(5,24)};
  \addlegendentry{$5f{-}1$ — \fl{quorum.twostep}}
\end{dgdata}
}

%%blk intro.notonce.06 from=c2336f35c2a2
Η ίδια λογική διέπει τη δεύτερη μεγάλη διάσταση, την \emph{αλλαγή όψης}
(\en{view change}). Στη διαδρομή χωρίς σφάλματα τα πρωτόκολλα διαφέρουν
κυρίως ως προς την τοπολογία και το πλήθος των φάσεων. Αποκλίνουν
δραστικότερα στην αντικατάσταση του ηγέτη, που οι σχεδιαστές του HotStuff
ονομάζουν «\en{the epicenter of replication}» \cite{yin2019hotstuff}. Ο
Πίνακας~\ref{tab:viewchange} συνοψίζει τέσσερις χαρακτηριστικές θέσεις σε
αυτόν τον άξονα.

%%blk intro.notonce.tab.viewchange from=b62f864471a2
\begin{table}[t]
\centering
\caption[Τέσσερις χαρακτηριστικές θέσεις στον άξονα της αλλαγής όψης:
κόστος, αποκρισιμότητα και πλήθος φάσεων.]%
{Τέσσερις χαρακτηριστικές θέσεις στον άξονα της αλλαγής όψης:
κόστος, αποκρισιμότητα και πλήθος φάσεων. Ο πίνακας αναπαράγει τον
Πίνακα~1 του \cite{malkhi2023hotstuff2}\prep\ μαζί με την επιφύλαξη που
η ίδια πηγή διατυπώνει. Δεν πρέπει να διαβαστεί χωρίς αυτήν. Δεν είναι
εξαντλητικός, επειδή η ίδια πηγή απαριθμεί άλλες σχεδιάσεις στον ίδιο
άξονα.}
\label{tab:viewchange}
\small
\begin{tabular}{@{}l c c c@{}}
\toprule
\textbf{Πρωτόκολλο} & \textbf{Κόστος αλλαγής όψης} &
\textbf{Αποκρίσιμη (\en{responsive})} & \textbf{Φάσεις} \\
\midrule
PBFT \cite{castro1999pbft}        & $O(n^{2})$ & ναι  & 2 \\
Tendermint \cite{buchman2018gossip}\prep & $O(n)$   & όχι  & 2 \\
HotStuff \cite{yin2019hotstuff}   & $O(n)$     & ναι  & 3 \\
HotStuff-2 \cite{malkhi2023hotstuff2}\prep & $O(n)$ & ναι & 2 \\
\bottomrule
\end{tabular}

\vspace{4pt}
\begin{minipage}{0.94\linewidth}\setstretch{1.08}\footnotesize
\textbf{Επιφύλαξη της ίδιας της πηγής.} Η πηγή
\cite{malkhi2023hotstuff2}\prep\ σημειώνει ρητά ότι τα PBFT και
Tendermint έχουν κόστος αλλαγής όψης $O(n^{3})$ και $O(n^{2})$
αντίστοιχα, και ότι τα φράγματα του πίνακα προκύπτουν αφού και τα δύο
βελτιωθούν κατά παράγοντα $O(n)$ --- βελτίωση που η ίδια πηγή αποδίδει
αλλού στη συνάθροιση υπογραφών. Η αναλυτική μέτρηση πολυπλοκότητας στο
\cite{yin2019hotstuff} (Πίνακας~1) δίνει για το PBFT $O(n^{2})$
αυθεντικοποιητές με σταθερό ηγέτη, $O(n^{3})$ σε αλλαγή όψης και
$O(f n^{3})$ σε αλληλουχία $f$ αποτυχιών ηγέτη. Δεν πρόκειται για
ανεξάρτητη επιβεβαίωση, καθώς οι δύο εργασίες μοιράζονται συγγραφέα. Το
\cite{jalalzai2023fasthotstuff} επιστρέφει σε δύο γύρους διατηρώντας την
αποκρισιμότητα, αλλά δηλώνει ρητά ότι \emph{δεν} αποφεύγει τη μεταφορά
τετραγωνικού πλήθους μηνυμάτων αλλαγής όψης όταν ο ηγέτης αποτύχει. Αυτό
που μειώνει είναι το πλήθος των υπογραφών που επαληθεύει κάθε κόμβος, από
$O(n)$ σε δύο. Οι Malkhi και Nayak δημοσιεύουν εκτεταμένη περίληψη σε
ePrint, που δεν έχει περάσει από κρίση \cite{malkhi2023hotstuff2}\prep.
Αναφέρεται ως τέτοια.

\smallskip
\textbf{Φάση δεν σημαίνει καθυστέρηση μηνύματος.} Ο πίνακας μετρά
\emph{φάσεις}, με τη σύμβαση της πηγής του. Μια φάση του HotStuff
κοστίζει δύο διαδοχικές καθυστερήσεις, επειδή περνά από τον ηγέτη. Μια
φάση ψηφοφορίας του PBFT κοστίζει μία, επειδή είναι όλοι-προς-όλους. Οι
δύο μονάδες δεν συγκρίνονται, και η διάκριση ορίζεται μία φορά στην
Ενότητα~\ref{sec:performance}.
\end{minipage}
\end{table}

%%blk intro.notonce.07 from=21b30534c3c2
Ο πίνακας δείχνει ότι κανένα από τα τέσσερα πρωτόκολλα δεν είναι
«καλύτερο». Το καθένα πληρώνει διαφορετικό λογαριασμό. Το Tendermint
καταργεί το υποπρωτόκολλο αλλαγής όψης και πληρώνει με αναμονή του
χειρότερου δυνατού $\Delta$. Δηλαδή πληρώνει με την \emph{αποκρισιμότητα}
(\en{responsiveness}), την ιδιότητα να προχωρά το πρωτόκολλο με τον
πραγματικό και όχι με τον μέγιστο χρόνο του δικτύου
\cite{buchman2018gossip}\prep\ και \cite{malkhi2023hotstuff2}\prep. Το
HotStuff πληρώνει με τρίτη φάση \cite{yin2019hotstuff}. Το HotStuff-2
ισχυρίζεται ότι κανένα από τα δύο τιμήματα δεν ήταν αναγκαίο
\cite{malkhi2023hotstuff2}\prep.

%%blk intro.opengap.head from=9649f2a47d7e
\section{Γιατί η σύγκριση παραμένει ανοιχτό ερώτημα}
\label{sec:opengap}

%%blk intro.opengap.01 from=9d0602b13d5f
Ο χώρος σχεδίασης είναι γνωστός και οι πολυπλοκότητες υπολογισμένες. Θα
μπορούσε επομένως κανείς να υποθέσει ότι η σύγκριση των πρωτοκόλλων είναι
λυμένο πρόβλημα. Δεν είναι, για τρεις λόγους που τεκμηριώνονται από την
ίδια τη βιβλιογραφία.

%%blk intro.opengap.02 from=2495649c056e
\paragraph{Πρώτον: τα δημοσιευμένα νούμερα δεν μεταφέρονται.}
Το Σχήμα~\ref{fig:scaling} το δείχνει στη δυσκολότερη δυνατή περίπτωση
για την ένσταση: ίδια εργασία, ίδια ενιαία υλοποίηση, ίδια μετρική, ίδιο
εύρος $n$. Αλλάζει μόνο η τοπολογία ανάπτυξης. Το κόστος κλιμάκωσης του
PBFT περνά από 65\% σε 84\%, του PoE από 39\% σε 51\% και του FTB από
32\% σε 36\% \cite{amiri2024bedrock}. Η ίδια εργασία δείχνει και
αντιστροφή κατάταξης χωρίς καμία αλλαγή περιβάλλοντος. Το SBFT αποδίδει
12\% \emph{υψηλότερο} \en{throughput} από το HotStuff σε $n=4$ και 8\%
\emph{χαμηλότερο} σε $n=100$, στο ίδιο πείραμα \cite{amiri2024bedrock}.

%%blk intro.opengap.03 from=780afdd88510
Όταν αλλάξει και η υλοποίηση, οι διαφορές γίνονται τάξεις μεγέθους. Σε
προσομοίωση πλανητικού δικτύου με $n=256$ κόμβους, το Kauri αποδίδει
περίπου εικοσαπλάσιο \en{throughput} από το HotStuff-SECP256K1, δηλαδή
4.917 πράξεις ανά δευτερόλεπτο έναντι 246. Εδώ κάθε πρωτόκολλο έχει
διαφορετική υλοποίηση, και οι τιμές είναι προσομοιωμένες και όχι
μετρημένες \cite{berger2023simulation}.

%%blk intro.opengap.fig.scaling from=bca3f3287031
\dgfigure{fig:scaling}%
{Το ίδιο πείραμα σε δύο τοπολογίες ανάπτυξης: ποσοστιαία μείωση του
\protect\en{throughput} καθώς οι κόμβοι αυξάνονται από 4 σε 100.}%
{Το ίδιο πείραμα σε δύο τοπολογίες ανάπτυξης: ποσοστιαία μείωση του
\protect\en{throughput} καθώς οι κόμβοι αυξάνονται από 4 σε 100. Όλα τα
νούμερα προέρχονται από την ίδια εργασία και την ίδια ενιαία υλοποίηση
των πρωτοκόλλων \cite{amiri2024bedrock}. \textbf{(α)}~Ανάπτυξη σε ένα
\en{datacenter}: Amazon EC2, στιγμιότυπα \texttt{c4.2xlarge} (8 vCPU,
15\,GB RAM, Intel Xeon E5-2666~v3 στα 3{,}50\,GHz), ένας κόμβος ανά
εικονική μηχανή, μέγεθος παρτίδας 400, ωφέλιμο φορτίο 128/128\,\en{byte},
πελάτες σε κλειστό βρόχο \cite{schroeder2006openclosed}\unver, μέσος όρος πέντε
εκτελέσεων, χωρίς σφάλματα. Τρεις διευκρινίσεις για το σχήμα. Το FaB
απαιτεί $5f{+}1$, άρα ξεκινά από $n=6$. Το \cite{amiri2024bedrock}
υλοποιεί το HotStuff αλυσιδωτά, με \en{pipelining}. Τα FTB και FLB είναι
νέα πρωτόκολλα που εισάγει η ίδια εργασία, όχι καθιερωμένα.
\textbf{(β)}~Τα τρία πρωτόκολλα για τα οποία η ίδια εργασία δίνει και
γεωκατανεμημένο νούμερο (Παράρτ.~D, ίδιο πλαίσιο υλοποίησης, ίδια
μετρική, ίδια παρτίδα 400, χωρίς σφάλματα). \emph{Δύο επιφυλάξεις που
οφείλει να δηλώσει η ίδια η εργασία.} Στο γεωκατανεμημένο σύνολο το
\cite{amiri2024bedrock} αυξάνει το \en{pipelining stretch} των Kauri και
FTB από 4 σε 6, και το εύρος του FTB φτάνει στο $n=99$ αντί για $n=100$.
Καμία από τις δύο δεν αφορά το PoE, άρα οι καθαρές --- χωρίς καμία άλλη
μεταβολή --- συγκρίσεις είναι δύο: του \textbf{PBFT} (65\,\% $\to$
84\,\%) και του \textbf{PoE} (39\,\% $\to$ 51\,\%). Και για τα τρία,
πάντως, ο αλγόριθμος και η μετρική παραμένουν ίδιοι, και αυτό που αλλάζει
είναι η τοπολογία ανάπτυξης.}%
{% fig-scaling — the same experiment, two deployment topologies. Two panels, so
% two dgbarchart environments stacked: the standard fixes every figure at
% \dgwidth, which leaves no room to set them side by side.
% All seven protocols are Byzantine, so the LAN series is dgBFT in both panels.
% Panel (b)'s second series is a second DEPLOYMENT of the same Byzantine
% protocols, so it takes a further-series colour (dgAux) — dgCFT would claim
% these are crash-fault protocols.
\begin{dgbarchart}[
  title={\fl{scaling.panel.lan}},
  ylabel={\fl{scaling.ylabel} (\%)},
  ymax=78,
  symbolic x coords={PBFT,FaB,SBFT,HotStuff,PoE,FTB,Kauri}, xtick=data,
  nodes near coords, every node near coord/.append style={dgtag, fill=dgPaper},
]
\addplot[dgBFT, fill=dgBFT] coordinates
  {(PBFT,65) (FaB,63) (SBFT,55) (HotStuff,45) (PoE,39) (FTB,32) (Kauri,31)};
\end{dgbarchart}

\par\vspace{5mm}

\begin{dgbarchart}[
  title={\fl{scaling.panel.geo}},
  ylabel={\fl{scaling.ylabel} (\%)},
  ymax=104,
  symbolic x coords={PBFT,PoE,FTB}, xtick=data,
  legend pos=north east,
  nodes near coords, every node near coord/.append style={dgtag, fill=dgPaper},
]
\addplot[dgBFT, fill=dgBFT,
         every node near coord/.append style={xshift=-2pt}]
  coordinates {(PBFT,65) (PoE,39) (FTB,32)};
  \addlegendentry{LAN}
\addplot[dgAux, fill=dgAux,
         every node near coord/.append style={xshift=2pt}]
  coordinates {(PBFT,84) (PoE,51) (FTB,36)};
  \addlegendentry{\fl{scaling.legend.geo}}
\end{dgbarchart}
}

%%blk intro.opengap.04 from=9a4a2e527cee
\paragraph{Δεύτερον: η ίδια σύγκριση αντιστρέφεται ανάλογα με το τι
μετριέται.}
Το πιο διδακτικό παράδειγμα είναι το EPaxos. Η πρωτότυπη εργασία
ισχυρίζεται βέλτιστη καθυστέρηση δέσμευσης σε δίκτυο ευρείας περιοχής,
ομοιόμορφη κατανομή φορτίου και ομαλή υποβάθμιση, με πειράματα σε Amazon
EC2 \cite{moraru2013epaxos}. Η ανεξάρτητη επαναξιολόγηση των Tollman
κ.ά.\ καταλήγει σε τρία συμπεράσματα \cite{tollman2021epaxosrevisited}.
\emph{Επιβεβαιώνει} τον ισχυρισμό για τη \emph{διάμεση} καθυστέρηση.
Δείχνει όμως καθυστέρηση ουράς (\en{tail latency}) πάνω από τέσσερις
φορές χειρότερη από το Multi-Paxos. Δείχνει επίσης ότι το ποσοστό των
πράξεων που ολοκληρώνονται σε έναν κύκλο μετ' επιστροφής έπεσε στο 65\%
σε ένα από τα πειράματά τους, υπό φόρτο με υψηλή σύγκρουση εντολών τον
οποίο η πρωτότυπη μεθοδολογία δεν προκαλούσε. Οι Matte κ.ά.\ αλλάζουν
πάλι τη μετρική, από απόλυτο \en{throughput} σε αποδοτικότητα πόρων.
Καταλήγουν ότι σε περιβάλλον σταθερού προϋπολογισμού πόρων, τυπικό των
πραγματικών νεφών, το μη βελτιστοποιημένο Multi-Paxos επιτυγχάνει
\emph{σχεδόν} διπλάσιο \en{throughput} από το EPaxos
\cite{matte2021wasteful}. Και οι τρεις εργασίες είναι σωστές. Απλώς δεν
μετρούν το ίδιο πράγμα. Η συνέπεια για κάθε συγκριτική αποτίμηση είναι
ότι η επιλογή της μετρικής δεν είναι τεχνική λεπτομέρεια αλλά μέρος του
συμπεράσματος, και πρέπει να δηλώνεται πριν από τη μέτρηση.

%%blk intro.opengap.05 from=84898cb359c8
\paragraph{Τρίτον: η διαφορά του πρωτοκόλλου δεν βρίσκεται εκεί που την
ψάχνει η μέτρηση.}
Οι Howard και Mortier αναλύουν τα δύο πρωτόκολλα που κυριαρχούν στα
παραγωγικά συστήματα. Καταλήγουν ότι Paxos και Raft ακολουθούν «πολύ
παρόμοια προσέγγιση, διαφέροντας μόνο ως προς την εκλογή ηγέτη»
\cite{howard2020paxosraft}. Αν αυτό ισχύει, τότε μια μέτρηση που
συγκρίνει δύο \emph{συστήματα} τα οποία υλοποιούν αυτά τα πρωτόκολλα
μετράει, στη διαδρομή χωρίς σφάλματα, μηχανική υλοποίησης και όχι
αλγοριθμική διαφορά. Μια τέτοια μέτρηση επομένως \emph{απομονώνει} την
υλοποίηση, που είναι ακριβώς το σωστό πείραμα αν αυτό είναι το ερώτημα.
Στη διαδρομή αλλαγής ηγέτη ισχύει το αντίστροφο. Εκεί συγκεντρώνεται όλη
η αλγοριθμική διαφορά, και οι ίδιοι οι συγγραφείς κάνουν δύο χωριστές
δηλώσεις. Η πρώτη είναι δομική διαπίστωση: η φάση εκλογής του Raft είναι
\emph{ελαφρύτερη σε μηνύματα}, επειδή δεν μεταφέρει εγγραφές του
\en{log}. Η δεύτερη είναι ρητή πρόβλεψη: η εκλογή αναμένεται \emph{και
πιο αργή και με μεγαλύτερη διακύμανση} στον χρόνο
\cite{howard2020paxosraft}. Δύο επιφυλάξεις είναι υποχρεωτικές. Η
σύγκριση γίνεται πάνω σε μια \emph{απλοποιημένη} διατύπωση του Paxos στην
ορολογία του Raft, και οι συγγραφείς δηλώνουν ότι και οι δύο αλγόριθμοι
που παρουσιάζουν είναι αφελείς εκ σχεδιασμού \cite{howard2020paxosraft}.

%%blk intro.opengap.06 from=ce6b101bc4ee
Το κενό, τέλος, έχει διατυπωθεί ρητά, και έχει καλυφθεί μερικώς. Οι
Ailijiang κ.ά.\ διαπιστώνουν ότι τα πρωτόκολλα ισχυρής συνέπειας
«εμφανίζουν ευρέως μεταβαλλόμενες επιδόσεις υπό διαφορετικές συνθήκες
δικτύου, φόρτου και ανάπτυξης, και καμία προηγούμενη μελέτη δεν προσφέρει
ολοκληρωμένη ανατομία και σύγκριση των επιδόσεών τους» (μετάφραση δική
μας), και το καλύπτουν για την οικογένεια κατάρρευσης
\cite{ailijiang2019dissecting}. Οι Amiri κ.ά.\ κάνουν το αντίστοιχο για
τη βυζαντινή \cite{amiri2024bedrock}. Οι Alqahtani και Demirbas μετρούν
πρωτόκολλα και των δύο οικογενειών στο ίδιο πλαίσιο
\cite{alqahtani2021bottlenecks}\unver.

%%blk intro.opengap.07 from=4c5e67029ab6
Αυτό που παραμένει ανοιχτό δεν είναι επομένως η ενιαία μέτρηση καθαυτή.
Είναι η μέτρηση των \emph{παραγωγικών υλοποιήσεων} --- αυτών που
πραγματικά αναπτύσσονται --- με ενιαίο όργανο, ενιαίο μοντέλο φόρτου και
σενάρια σφάλματος. Και οι τρεις παραπάνω εργασίες μετρούν ερευνητικά
πρωτότυπα, ή ενιαίες επανυλοποιήσεις μέσα σε ένα πλαίσιο, και όχι τα
συστήματα που τρέχουν σε παραγωγή. Όπως δείχνει το
Σχήμα~\ref{fig:scaling}, η επανυλοποίηση είναι ακριβώς η μεταβλητή που
μπορεί να μετακινήσει το αποτέλεσμα.

%%blk intro.aims.head from=ce135bb1be78
\section{Σκοπός, ερευνητικά ερωτήματα και όρια της εργασίας}
\label{sec:aims}

%%blk intro.aims.01 from=1d3ab27b5434
Σύμφωνα με την εγκεκριμένη πρόταση εκπόνησης (βλ.\ Παράρτημα~Α), η
εργασία έχει τρεις σκοπούς. Ο πρώτος είναι η μελέτη της εξέλιξης των
πρωτοκόλλων συναίνεσης. Ο δεύτερος είναι η θεωρητική και αρχιτεκτονική
ανάλυση των κυριότερων από αυτά. Ο τρίτος είναι η σύγκριση των επιδόσεων
επιλεγμένων υλοποιήσεων σε αντιπροσωπευτικά σενάρια. Τα κριτήρια της
σύγκρισης είναι τέσσερα: η καθυστέρηση συναίνεσης, το \en{throughput}, η
ανοχή σε σφάλματα και η κλιμάκωση. Τα ερευνητικά ερωτήματα που
ακολουθούν εξειδικεύουν τους σκοπούς αυτούς.

%%blk intro.aims.rq from=5fdf024d64a8
\begin{description}[leftmargin=2.1em,style=nextline,font=\normalfont\bfseries]
  \item[ΕΕ1 --- Εξέλιξη.] Ποιες συγκεκριμένες ανταλλαγές (\en{trade-offs})
  εξηγούν κάθε μετακίνηση μέσα στον χώρο σχεδίασης, και ποιο τίμημα
  πληρώνεται σε κάθε μία; Απαντάται αναλυτικά, από τις πρωτογενείς πηγές.
  \item[ΕΕ2 --- Κοινή δομή.] Ποια αναλλοίωτα (\en{invariants}) μοιράζονται
  όλα τα πρωτόκολλα, ώστε να μπορούν να περιγραφούν σε ενιαίο σχήμα αντί
  για κατάλογο ανεξάρτητων αλγορίθμων;
  \item[ΕΕ3 --- Συγκρισιμότητα.] Υπό ποιες προϋποθέσεις είναι θεμιτή η
  απευθείας σύγκριση δύο υλοποιήσεων, και πότε η μέτρηση αποτυπώνει
  μηχανική υλοποίησης αντί για ιδιότητα του πρωτοκόλλου;
  \item[ΕΕ4 --- Εμπειρική συμπεριφορά.] Πώς συμπεριφέρονται επιλεγμένες
  παραγωγικές υλοποιήσεις σε ελεγχόμενο περιβάλλον μικρής συστάδας ως
  προς καθυστέρηση, \en{throughput} και ανοχή σε σφάλματα; Η εργασία
  δίνει δικές της μετρήσεις γι' αυτό στο κεφάλαιο της αποτίμησης.
  \item[ΕΕ5 --- Κλιμάκωση.] Πώς μεταβάλλεται η συμπεριφορά καθώς μεγαλώνει
  το \en{quorum}; Η εργασία προσεγγίζει το ερώτημα σε δύο επίπεδα, και τα
  δηλώνει χωριστά. Το πειραματικό επίπεδο καλύπτει μόνο μικρά μεγέθη
  συστάδας, όπου ελέγχεται το κόστος της αριθμητικής του
  Σχήματος~\ref{fig:quorum}. Το βιβλιογραφικό επίπεδο καλύπτει τα μεγάλα
  μεγέθη, όπου εμφανίζεται η ασυμπτωτική συμπεριφορά. Η εργασία δίνει
  δικές της μετρήσεις για το πειραματικό επίπεδο στο κεφάλαιο της
  αποτίμησης.
\end{description}

%%blk intro.aims.02 from=dfce2af38acb
Η εργασία δηλώνει τα όρια εδώ και όχι στα συμπεράσματα. Πρώτον, η εργασία
αποτιμά \emph{υλοποιήσεις} και όχι πρωτόκολλα. Το μετρήσιμο αντικείμενο
είναι πάντοτε ένα σύστημα με συγκεκριμένη έκδοση και συγκεκριμένη
ρύθμιση. Κάθε ισχυρισμός για το υποκείμενο πρωτόκολλο περνά υποχρεωτικά
από το ΕΕ3.

%%blk intro.aims.03 from=d606b8a15bc8
Δεύτερον, ο διαθέσιμος προϋπολογισμός υποδομής περιορίζει το εύρος των
μεγεθών συστάδας που μπορεί να μελετήσει ένα πείραμα. Το εύρος αυτό είναι
μικρό σε σχέση με τις κλίμακες όπου εμφανίζονται οι ενδιαφέρουσες
ασυμπτωτικές συμπεριφορές. Αυτός είναι ο λόγος του διαχωρισμού του ΕΕ5 σε
δύο επίπεδα.

%%blk intro.aims.04 from=67089110cff9
Τρίτον, ένα πρωτόκολλο μπορεί να καλυφθεί \emph{αναλυτικά} χωρίς δικές
μας μετρήσεις. Όπου συμβαίνει αυτό, το κεφάλαιο το δηλώνει στην αρχή
του. Οι όποιοι αριθμοί προέρχονται εκεί από τρίτες δημοσιεύσεις που
κατονομάζονται ρητά, και δεν στέκουν ποτέ στον ίδιο πίνακα με δικές μας
μετρήσεις. Το πλήρες πλαίσιο εγκυρότητας, τα μεθοδολογικά κριτήρια και οι
απειλές εγκυρότητας ανήκουν στο κεφάλαιο της μεθοδολογίας. Το παρόν
κεφάλαιο δεν τα προεξοφλεί.

%%blk intro.structure.head from=330c145eb75d
\section{Συνεισφορά και δομή της εργασίας}
\label{sec:structure}

%%blk intro.structure.01 from=241b572f4703
Η εργασία επιδιώκει τέσσερις συνεισφορές.

%%blk intro.structure.list from=4b7e7f203919
\begin{enumerate}[label=\textbf{(\alph*)},leftmargin=2.4em]
  \item Μια χρονολογημένη αφήγηση της εξέλιξης των πρωτοκόλλων, με ρητό
  κανόνα ταξινόμησης. Η αφήγηση αποδίδει κάθε μετάβαση σε κατονομασμένη
  ανταλλαγή.
  \item Μια ενιαία δομική περιγραφή. Το ίδιο σχήμα ανάγνωσης καλύπτει τα
  πρωτόκολλα και των δύο μοντέλων σφαλμάτων.
  \item Μια συγκριτική αποτίμηση επιλεγμένων παραγωγικών υλοποιήσεων σε
  ελεγχόμενο περιβάλλον, με προβλέψεις δηλωμένες εκ των προτέρων
  (\en{pre-registered predictions}). Η εργασία δηλώνει ρητά και τι
  \emph{δεν} μπορεί να συναχθεί από τα δεδομένα. Η εργασία δίνει δικές
  της μετρήσεις γι' αυτό στο κεφάλαιο της αποτίμησης.
  \item Ένα πλαίσιο απόφασης που αντιστοιχίζει απαιτήσεις σε οικογένειες
  πρωτοκόλλων.
\end{enumerate}

%%blk intro.structure.02 from=ec1f8fc7af8b
Η διάρθρωση ακολουθεί δύο χρονολογικές διαδρομές, μία ανά μοντέλο
σφαλμάτων, αντί για μία αυστηρά ενιαία. Η επιλογή είναι συνειδητή
απόκλιση από την αυστηρά χρονολογική παρουσίαση και τίθεται ρητά υπόψη
της επιβλέπουσας. Με ενιαία σειρά, η ακολουθία του
Σχήματος~\ref{fig:timeline} υποχρεώνει τον αναγνώστη να εναλλάξει μοντέλο
σφαλμάτων πέντε φορές. Το αποτέλεσμα είναι να παρατίθενται δίπλα-δίπλα
μεγέθη που δεν συγκρίνονται.

%%blk intro.structure.03 from=1bfe64225865
Το Κεφάλαιο~\ref{ch:background} θεμελιώνει το θεωρητικό υπόβαθρο: μοντέλα
σφαλμάτων και συγχρονισμού, SMR, \en{quorums} και κριτήρια ορθότητας. Τα
κεφάλαια της πρώτης διαδρομής καλύπτουν την οικογένεια κατάρρευσης, και
τα κεφάλαια της δεύτερης τη βυζαντινή. Κάθε πρωτόκολλο ακολουθεί κοινό
πρότυπο πέντε υποενοτήτων:

%%blk intro.structure.template from=14f76f86d68c
\begin{enumerate}[label=\arabic*.,leftmargin=2.4em]
  \item ιστορικό πλαίσιο και πρωτογενής πηγή·
  \item ο αλγόριθμος·
  \item η υλοποίηση που αποτιμήθηκε, ή η αιτιολόγηση της απουσίας της·
  \item οι μετρήσεις·
  \item η ανάγνωση των αποτελεσμάτων έναντι της πρόβλεψης που είχε δηλωθεί
  εκ των προτέρων.
\end{enumerate}

%%blk intro.structure.04 from=162cf9484024
Ακολουθεί το κεφάλαιο της μεθοδολογίας και της εγκυρότητας. Έπειτα
έρχονται η συγκεντρωτική σύγκριση, το πλαίσιο απόφασης και, τέλος, τα
συμπεράσματα με τους περιορισμούς και τη μελλοντική εργασία.

%%blk intro.sources.head from=d131226a5c4a
\section{Σημείωση για την ποιότητα των πηγών}
\label{sec:sources}

%%blk intro.sources.01 from=db5ff35becc9
Μεγάλο μέρος της βιβλιογραφίας των πρωτοκόλλων συναίνεσης κυκλοφορεί ως
\en{preprint}, ως \en{whitepaper} ή ως τεχνική τεκμηρίωση έργου. Γι' αυτό
τον λόγο η εργασία εφαρμόζει τρεις κανόνες σε όλα της τα κεφάλαια.
Πρώτον, κάθε αλγόριθμος περιγράφεται από την πρωτογενή του πηγή. Οι
ανασκοπήσεις και οι ενιαίες πλατφόρμες χρησιμοποιούνται για ενοποιητικά
σχήματα και για ταξινομήσεις. Όπου μια τέτοια εργασία μεταφέρει φράγμα ή
αριθμό τρίτου, πράγμα που συμβαίνει μόνο κατ' εξαίρεση, κατονομάζεται και
η πρωτογενής πηγή, όπως στη λεζάντα του Σχήματος~\ref{fig:quorum}.

%%blk intro.sources.02 from=d0da0f8f2fc7
Δεύτερον, οι μη κριθείσες πηγές δηλώνονται ως τέτοιες στο σημείο όπου
χρησιμοποιούνται. Η σήμανση \grey\ δείχνει τη γκρίζα βιβλιογραφία:
\en{whitepapers}, KIP και ιστοσελίδες. Η σήμανση \prep\ δείχνει τα
\en{preprints} και τις εκτεταμένες περιλήψεις χωρίς κρίση. Μια τρίτη
σήμανση, \unver, λέει κάτι διαφορετικό από τις δύο προηγούμενες: λέει ότι
τα βιβλιογραφικά στοιχεία της ίδιας της πηγής δεν ήταν δυνατόν να
ελεγχθούν ούτε στο άρθρο ούτε σε ευρετήριο, και πάρθηκαν από τον κατάλογο
αναφορών άλλης εργασίας. Ο αναγνώστης που θέλει να ακολουθήσει την
παραπομπή πρέπει να το ξέρει πριν το επιχειρήσει. Τρίτον, κάθε
δημοσιευμένος αριθμός που αναπαράγει η εργασία συνοδεύεται από το
περιβάλλον στο οποίο παρήχθη: υλικό, μοντέλο φόρτου, τοπολογία ανάπτυξης,
μέγεθος συστάδας και πλήθος επαναλήψεων. Οι λεζάντες των
Σχημάτων~\ref{fig:quorum} και~\ref{fig:scaling} το κάνουν αυτό. Αριθμός
χωρίς περιβάλλον δεν είναι τεκμήριο.

%%blk intro.sources.03 from=1662413bd872
Η ίδια αρχή εφαρμόζεται και στα βιβλιογραφικά στοιχεία. Στο αρχείο
βιβλιογραφίας κάθε εγγραφή φέρει σήμανση για το πώς επαληθεύτηκε: από την
πρώτη σελίδα του ίδιου του άρθρου, από αυθεντικό ευρετήριο μέσω του DOI,
ή καθόλου. Στην τελευταία περίπτωση η εγγραφή δηλώνεται ως μη
επαληθευμένη. Η διάκριση δεν είναι τυπολατρική. Τρεις περιπτώσεις όπου η
σήμανση απέτρεψε λάθος αξίζει να καταγραφούν.

%%blk intro.sources.04 from=f76fb1b59cec
Η πρώτη περίπτωση είναι το \cite{berger2023simulation}. Το αντίτυπο που
συλλέχθηκε είναι \en{preprint} και θα είχε δηλωθεί ως τέτοιο, ενώ το έργο
έχει κριθείσα εκδοχή σε συνέδριο. Η σήμανση \emph{επαλήθευσης} περιγράφει
το αντίτυπο που κρατάμε. Οι σημάνσεις \prep\ και \grey\ περιγράφουν το
\emph{έργο}. Εδώ τα δύο αποκλίνουν, και η απόκλιση εντοπίζεται μόνο με
έλεγχο ευρετηρίου. Το αντίτυπο είναι \en{preprint}, το έργο όχι, άρα το
\cite{berger2023simulation} \emph{δεν} φέρει \prep.

%%blk intro.sources.05 from=3e325a7e9e8a
Η δεύτερη περίπτωση είναι το \cite{oki1988viewstamped}. Καταχωρείται στο
Crossref με διαφορετική διατύπωση τίτλου από εκείνη των πρακτικών, και η
πλήρης μορφή επιβεβαιώνεται από τη βιβλιογραφία του ίδιου του
\cite{lamport1998parttime}. Η τρίτη περίπτωση είναι το
\cite{jalalzai2023fasthotstuff}. Ο τίτλος του είναι «\en{\dots A Fast and
Robust BFT Protocol}», όχι «\en{Resilient}», όπως εμφανίζεται σε
ορισμένους καταλόγους.

%%blk intro.sources.06 from=8fdfc7517d46
Μία τέταρτη περίπτωση δείχνει το ίδιο σε επίπεδο χαρακτήρα. Το
\en{corpus} έχει \emph{δύο} αντίτυπα του \cite{lamport1998parttime}, την
έκδοση της ACM και τη διορθωμένη του συγγραφέα. Τυπώνουν την ίδια
πληροφορία με άλλα γράμματα: «\en{Received: January 1990; accepted: March
1998}» η πρώτη, «\en{Received January 1990; Accepted March 1998}» η
δεύτερη. Όποιο και να παρατεθεί ως αυτούσιο, θα είναι λανθασμένο για το
άλλο. Γι' αυτό η λεζάντα του Σχήματος~\ref{fig:timeline} δίνει το
\emph{περιεχόμενο}. Όταν δύο αντίτυπα διαφέρουν, το αυτούσιο παράθεμα
απαιτεί να δηλωθεί και το αντίτυπο.
